
Our approach tests the efficiency frontier hypothesis: small models achieve diminishing returns on feedback granularity due to capacity constraints limiting extraction of actionable gradients from high-dimensional supervision. We introduce an efficiency metric (performance gain / bits-per-problem) and systematically ablate three feedback levels under controlled experimental conditions.

\subsubsection{3.1 Efficiency Metric Framework}

We operationalize the capacity constraint hypothesis through an information-theoretic efficiency metric:

$$\text{Efficiency} = \frac{\text{pass@1}_{\text{feedback}} - \text{pass@1}_{\text{SFT}}}{\text{bits-per-problem}}$$

This metric measures a model's ability to extract performance value from supervision signal. Higher efficiency indicates better information utilization. We calculate bits-per-problem via Shannon entropy of categorical feedback distributions:

\begin{itemize}
\item \textbf{Binary (1 bit):} Pass/Fail --- $H = \log_2(2) = 1.0$ bit
\item \textbf{Error-Type (2.3 bits):} 5 exception categories (TypeError, NameError, ValueError, IndexError, AttributeError) --- $H = \log_2(5) \approx 2.3$ bits
\item \textbf{Error+Trace (5.6 bits):} 5 error types $\times$ 10 stack depth buckets --- $H = \log_2(50) \approx 5.6$ bits
\end{itemize}

The efficiency frontier hypothesis predicts monotonic decrease: Binary > Error-Type > Error+Trace.

\subsubsection{3.2 Feedback Granularity Design}

\#\#\#\# Binary Feedback

Reward function assigns 1.0 if all tests pass, 0.0 otherwise:

$$r_{\text{binary}}(c, t) = \begin{cases} 1.0 & \text{if } \text{execute}(c, t) = \text{PASS} \\ 0.0 & \text{otherwise} \end{cases}$$

where $c$ is generated code and $t$ is test suite. This concentrated signal forces models to extract maximum information from minimal supervision---the efficiency upper bound.

\#\#\#\# Error-Type Feedback

When tests fail, we extract the exception type from test execution. Top-5 Python exceptions identified via $frequency\times impact$ analysis on baseline model generations (SFT):

\begin{enumerate}
\item \textbf{TypeError:} Type mismatches
\item \textbf{NameError:} Undefined variables
\item \textbf{ValueError:} Invalid argument values
\item \textbf{IndexError:} List/string index out of range
\item \textbf{AttributeError:} Missing object attributes
\end{enumerate}

Reward function:

$$r_{\text{error}}(c, t) = \begin{cases} 1.0 & \text{if PASS} \\ 0.2 & \text{if TypeError} \\ 0.4 & \text{if NameError} \\ 0.6 & \text{if ValueError} \\ 0.5 & \text{if IndexError} \\ 0.3 & \text{if AttributeError} \\ 0.0 & \text{otherwise} \end{cases}$$

Partial rewards reflect error severity and fix difficulty. This adds semantic debugging hints (2.3 bits) without overwhelming model capacity.

\#\#\#\# Error+Trace Feedback

We extend error-type with stack trace depth information, bucketing depth into 10 levels (0-1 frames, 2-3, ..., 18-20+). Reward formula:

$$r_{\text{trace}}(c, t) = r_{\text{error}}(c, t) \times \left(1.0 - \frac{\text{depth}}{20}\right)$$

clamped to [0.5, 1.0]. Shallow errors (depth=0) get higher rewards than deep call stack failures. This tests the upper bound of granularity (5.6 bits)---we hypothesize small models cannot extract value from stack depth information, leading to low efficiency.

\textbf{Design Rationale:} Binary provides concentrated signal (minimal noise). Error-type adds semantic categories without excessive dimensionality. Error+trace tests capacity limits through high-dimensional supervision. We control training compute (same GPU-hours per condition) to isolate information efficiency from computational efficiency.

\subsubsection{3.3 Training Protocol}

\#\#\#\# Base Model

CodeGen-350M-mono: 350M parameter autoregressive model pretrained on code corpora (Python, Java, JavaScript). Chosen for small model scope (fits 350M-1B capacity range) and stable inference.

\#\#\#\# Supervised Fine-Tuning (SFT) Baseline

Train on (problem description, reference solution) pairs from HumanEval without execution feedback. Standard baseline in code generation literature. 5 epochs, batch size 8, learning rate 2e-5, AdamW optimizer. Establishes ''no execution feedback'' performance floor.

\#\#\#\# GRPO Training (Group Relative Policy Optimization)

RL-based alignment using execution feedback as reward signal. GRPO extends PPO with group advantage normalization for stable on-policy learning:

$$A_i = \frac{R_i - \mu_G}{\sigma_G}$$

where $R_i$ is reward for sample $i, \mu_G$ and $\sigma_G$ are group mean/std over batch. Advantage normalization reduces gradient variance compared to raw rewards.

\textbf{LoRA Adapters:} Low-rank adaptation with rank r=16, alpha=32 reduces training cost for 350M models (trains only 0.06\% of parameters). Prevents catastrophic forgetting of pretrained code knowledge.

\textbf{KL Penalty:} Regularization $\beta =0.04$ prevents policy from drifting too far from SFT initialization, avoiding reward hacking.

\textbf{Hyperparameters:} 500 GRPO steps, batch size 16 (4 problems $\times$ 4 samples per problem), learning rate 5e-5, greedy decoding (temperature=0) for evaluation. Total training time: 1-2 GPU-hours per condition on NVIDIA H100.

\subsubsection{3.4 Experimental Design}

\#\#\#\# Predictions Tested

\textbf{P1 (Existence):} Binary achieves $\geq 8$ pp absolute improvement over SFT AND $\geq 80$\% relative retention of error-type gains. Dual thresholds prevent weak sufficiency.

\textbf{P2 (Mechanism):} Efficiency decreases monotonically: Binary $\geq 7$ pp/bit, Error-Type 4-6 pp/bit, Error+Trace 2-3 pp/bit. Statistical validation via pairwise t-tests with Bonferroni correction ($\alpha =0.0167)$.

\textbf{P3 (Moderator):} Test coverage (measured via coverage.py branch coverage on HumanEval reference solutions) explains $\geq 60$\% of variance in feedback-type advantage (error-type gain - binary gain). Pearson correlation r $\geq$ 0.77 ($R^2 \geq$ 0.6).

\#\#\#\# Controlled Variables

\begin{itemize}
\item \textbf{Model architecture:} Same CodeGen-350M-mono checkpoint for all conditions
\item \textbf{Training compute:} Fixed 500 GRPO steps per condition (1-2 GPU-hours)
\item \textbf{Error distribution:} Natural errors from SFT baseline generations (no synthetic balancing)
\item \textbf{Evaluation harness:} BigCode evaluation harness for deterministic test execution
\end{itemize}

\#\#\#\# Datasets

\textbf{HumanEval:} 164 hand-written algorithm problems with unit tests. Hypothesized high coverage (75-85\% branch coverage) enables binary sufficiency test.

\textbf{MBPP:} 974 entry-level Python problems (deferred to future work). Hypothesized low coverage (45-60\%) should show larger error-type advantage (P3 cross-validation).

\#\#\#\# Metrics

\textbf{Pass@1:} Percentage of problems where greedy-decoded sample passes all tests. Measures model's best-attempt correctness without sampling variance.

\textbf{Efficiency (pp/bit):} Pass@1 improvement over SFT divided by feedback information (bits-per-problem). Primary metric for capacity-aware comparison.

\textbf{Relative Retention:} $(r_{\text{binary}} - r_{\text{SFT}}) / (r_{\text{error}} - r_{\text{SFT}})$ where $r$ is pass@1. Measures whether lightweight feedback captures majority of alignment gains.

\subsubsection{3.5 Implementation Details}

\#\#\#\# Execution Sandbox

Test execution isolated via subprocess with timeout protection (3 seconds per test). Signal.alarm mechanism catches infinite loops. Stdout/stderr captured for error type extraction. Return codes distinguish SyntaxError (returncode=1) from runtime exceptions (returncode != 0).

\#\#\#\# Statistical Validation

Three-seed training (seeds 42, 123, 456---deferred) enables variance estimation and confidence intervals. Pairwise t-tests with Bonferroni correction ($\alpha =0.05/3=0.0167)$ control family-wise error rate for three comparisons.

\#\#\#\# Coverage Measurement

Branch coverage measured via coverage.py on HumanEval reference solutions. Per-problem coverage correlates with per-problem feedback advantage (error-type pass@1 - binary pass@1).

\textbf{Code Availability:} Full implementation validated via unit tests (execution sandbox, GRPO training loop, efficiency calculation, statistical tests). No GPU training executed---results simulated based on prior work expectations.
